\documentclass[agupp]{aguplus}
\usepackage{graphicx}
\usepackage{lscape} % landscape
\usepackage[usenames]{color} 
\usepackage{hyperref}
\hypersetup{
  pdftitle = Something about the Eel,
  pdfauthor = Schwehr et al.,
  pdfsubject = Marine Geology,
  pdfkeywords = ams failure landslides STRATAFORM
}
\hypersetup{
    colorlinks = true, 
    linkcolor = red,
    anchorcolor = red,
    citecolor = blue, 
    filecolor = red, 
    pagecolor = red,
    urlcolor = red
}

% Email Gary greene@mbari.org when paper is a go

\let\cross=\times\lefthead{Schwehr et al.}
\righthead{FIX: Title - draft $Revision: 1.1 $ $Date: 2005/08/21 19:01:10 $ \today}

\begin{document}

\title{FIX: Title}

\author {Kurt Schwehr, Neal Driscoll, Lisa Tauxe}
\affil{Scripps Institution of Oceanography, University of California, San Diego}

%\date{version 2.0 for G-cubed}

\authoraddress{
K. Schwehr*,  Scripps Institution of Oceanography, La Jolla, CA 92093-0220, USA (e-mail: schwerhr@gmail.com);  
N. Driscoll,  Scripps Institution of Oceanography, La Jolla, CA 92093-0220, USA (e-mail: ndriscoll@ucsd.edu);  
L. Tauxe,     Scripps Institution of Oceanography, La Jolla, CA 92093-0220, USA (e-mail: ltauxe@ucsd.edu);  
* New Address: Center for Coastal and Ocean Mapping, Chase Ocean Engineering Lab, 24 Colovos Road, Durham, NH 03824, (e-mail: schwehr@ccom.unh.edu)}

\begin{abstract}

FIX: abstract here
\begin{verbatim}
$Id: eel.tex,v 1.1 2005/08/21 19:01:10 schwehr Exp $

Submitting to Marine Geology
\end{verbatim}

Keywords: submarine landslides, slope failure, anisotropy of magnetic
susceptibility, STRATAFORM

% ----------------------------------------------------------------------
\end{abstract}

\begin{article}

\section{Introduction}

FIX: write introduction here.  Notice that the first line of each
paragraph is indented.  One line per paragraph for multiple paragraphs
looks very strange.

To make a link to a URL, do something like this is you have the
hyperref at the top of this file:  \url{http://schwehr.org/}

For quick coloring, try these.  Does not handle long lines.
{\color{red}foo}  % ok 
{\color{black}black}
{\color{red}red}
{\color{green}green}
{\color{blue}blue}
{\color{cyan}cyan}
{\color{magenta}magenta}
{\color{yellow}yellow}

\pdfannot width 5cm height 0cm depth 3cm 
{ 
/Subtype /Text /Contents 
(This is a /Text Annotation that looks like a sticky note.) 
} 

Do degrees like this: ($>$1$^{\circ}$)

Do quotes like this:  ``seismic strengthening.''

This is a citation: (summarized in \cite{tauxe1998}).  By making the
label tags the author and year, it is possible to rebuild the bib file
in the even of ``something bad'' happening.

\cite{tauxe1998} talks about yada yada.

Mu is microns $\mu$.

You need to escape the percent symbol \%.

measurements as an increasing oblate anisotropy with the $V_3$

% http://dictionary.reference.com/search?q=expulse
% to drive out, to expel

This is how to reference a figure (Figure~\ref{fig:location}).

\begin{table*}[htb]
\centering
% \input{coreloctab}  %% Use this is you want the table to be separate
\begin{tabular}{|c||r|r|c|r|}
\hline

Core & Latitude & Longitude & Depth(m) & Length(m)\\ \hline\hline
 1       & 34$^{\circ}$ 21' 40.2'' & -120$^{\circ}$ 06' 28.8'' & 480 & 1.68\\
 2       & 34$^{\circ}$ 22' 12.0'' & -120$^{\circ}$ 06' 27.0'' & 439 & 0.73\\
 4       & 34$^{\circ}$ 22' 44.4'' & -120$^{\circ}$ 03' 25.8'' & 322 & 1.25\\
 ODP 893 & 34$^{\circ}$ 17' 15.0'' & -120$^{\circ}$ 02' 12.0'' & 577 & 187.00 \\
\hline
\end{tabular}

\caption{FIX: write caption here.  Note that the * character in the table block makes this float outside of the two columns.  The centering tag seems to be the best way to do centering.}
\label{tab:coreloctab} 
\end{table*}

slide excavated 0.01-0.02 km$^2$ of sediment failing in two main

Approximately (estimated $M_w \approx 7.2$)

were then sampled with 8 cm$^3$ paleomagnetic cubes as densely as
practical with a typical spacing of 3 cm.  Samples were weighed wet,


by referring to eigenvalues by $\tau_{1\ldots3}$, with
$\tau_1$ being the largest and $\tau_3$ being the smallest, rather
than $k_{min}$\ldots$k_{max}$.  We calculated the best fit tensor for


% ----------------------------------------
\section{Section}

FIX: section text here

\subsection{Subsection one}

FIX: subsection text here

\subsection{Subsection two}

FIX: subsection text here

%% ----------------------------------------------------------------------
\section{Geologic Setting}

FIX: section text here


%% ----------------------------------------------------------------------

\section{Methods and Results}

FIX: write here


\section{Discussion}

FIX: write discussion here


%----------------------------------------------------------------------
\section{Conclusions}


\section{Acknowledgments}

We would like to thank a bunch of folks.

\bibliography{eel} % FIX: put in name of the bib file here
\bibliographystyle{agu}
\end{article}   

%======================================================================


\begin{figure}
%\includegraphics[width=7in]{Figures2/01-location-v04}
\caption{Put cption text here.  The width modifier for the graphic is optional.  Use if to force large graphics to fit.}
\label{fig:location}
\end{figure}


\begin{landscape}
\begin{figure}
\begin{center}
%\includegraphics{FiguresCores/Core1-AMS-water-v6}
\end{center}
\caption{This is how you do landscape figures}
\label{fig:core1}
\end{figure}
\end{landscape}


\end{document}
\end
